\bta{2016}

\begin{enumerate}
	%\renewcommand{\labelenumi}{\arabic{enumi}.}
	% A(\Alph) a(\alph) I(\Roman) i(\roman) 1(\arabic)
	%设定全局标号series=example	%引用全局变量resume=example
	%[topsep=-0.3em,parsep=-0.3em,itemsep=-0.3em,partopsep=-0.3em]
	%可使用leftmargin调整列表环境左边的空白长度 [leftmargin=0em]
	\item
	一个质量为 $m$ 的粒子限制在半径为 $a$ 的圆周上运动。
	\begin{enumerate}
		%\renewcommand{\labelenumi}{\arabic{enumi}.}
		% A(\Alph) a(\alph) I(\Roman) i(\roman) 1(\arabic)
		%设定全局标号series=example	%引用全局变量resume=example
		%[topsep=-0.3em,parsep=-0.3em,itemsep=-0.3em,partopsep=-0.3em]
		%可使用leftmargin调整列表环境左边的空白长度 [leftmargin=0em]
		\item
		求相应的本征值和本征函数。
		\item 
		$t=0$ 时， $\Psi=A \sin ^{2} \theta $，其中$ \theta $为极角， $ A $为归一化因子，求 $\mathrm{t}$ 时刻的波函数 $\Psi(t)$。 
		\item 
		求 $t$ 时刻垂直于圆周角动量的平均值。
		
		
		
	\end{enumerate}


\banswer{
\begin{enumerate}
\item 用同一个角坐标 $\theta$，$H_0=-\hbar^2\partial_\theta^2/(2ma^2)$。
令 $\psi=e^{i\nu\theta}$，周期条件要求 $\nu=n\in\mathbb Z$，归一化得
\[
 \boxed{\phi_n(\theta)=\frac{e^{in\theta}}{\sqrt{2\pi}},\qquad
 E_n=\frac{n^2\hbar^2}{2ma^2}}.
\]
\item 由 $\int_0^{2\pi}\sin^4\theta\,d\theta=3\pi/4$，可取
$A=2/\sqrt{3\pi}$。再用
\[
 \sin^2\theta=\frac12-\frac14e^{2i\theta}-\frac14e^{-2i\theta},
\]
得 $c_0=\sqrt{2/3}$、$c_2=c_{-2}=-1/\sqrt6$，且三系数的模方之和为 $1$。
因此
\[
 \boxed{\Psi(\theta,t)=\sqrt{\frac23}\phi_0(\theta)
       -\frac1{\sqrt6}\bigl[\phi_2(\theta)+\phi_{-2}(\theta)\bigr]
         e^{-iE_2t/\hbar}}.
\]
\item 垂直圆面的角动量为 $L_z=-i\hbar\partial_\theta$，满足
$L_z\phi_n=n\hbar\phi_n$。能量表象也对角化 $L_z$，故
\[
 \boxed{\langle L_z\rangle_t=0\cdot\frac23
                   +2\hbar\cdot\frac16-2\hbar\cdot\frac16=0}.
\]
正反两向角动量等概率，随时间积累的相位不改变此结论。
\end{enumerate}
}

\item 
已知一质量为 $m$ 的粒子处于 $\delta$ 势阱中, $\mathrm{V}(x)=-V_{0} \delta(x),\left(V_{0}>0\right)$ 中运动。
\begin{enumerate}
	%\renewcommand{\labelenumi}{\arabic{enumi}.}
	% A(\Alph) a(\alph) I(\Roman) i(\roman) 1(\arabic)
	%设定全局标号series=example	%引用全局变量resume=example
	%[topsep=-0.3em,parsep=-0.3em,itemsep=-0.3em,partopsep=-0.3em]
	%可使用leftmargin调整列表环境左边的空白长度 [leftmargin=0em]
	\item
	求粒子的束缚态能量，以及归一化的束缚态波函数。
	\item 
	利用 Feynman-Hellmann 定理和 virial 定理, 求束缚态粒子的平均动能和平均势能。
	
	
	
\end{enumerate}

\banswer{
\begin{enumerate}
\item 束缚态 $E<0$，记 $\kappa=\sqrt{-2mE}/\hbar>0$。在 $x\ne0$ 处，
方程为 $\psi''=\kappa^2\psi$；平方可积性和原点连续性给出
\[
 \psi(x)=Ce^{-\kappa|x|}.
\]
将 Schr\"odinger 方程在原点两侧积分，得到导数跃变
\[
 \psi'(0^+)-\psi'(0^-)=-\frac{2mV_0}{\hbar^2}\psi(0).
\]
左边为 $-2\kappa C$，故 $\kappa=mV_0/\hbar^2$，从而
\[
 \boxed{E=-\frac{mV_0^2}{2\hbar^2}}.
\]
归一化 $1=2|C|^2\int_0^\infty e^{-2\kappa x}dx=|C|^2/\kappa$ 给出
\[
 \boxed{\psi(x)=\sqrt\kappa e^{-\kappa|x|}}.
\]
跃变条件唯一确定正数 $\kappa$，且奇态在原点为零无法产生束缚，故只有此一个束缚态。
\item 用 $V_0$ 为参数，由 Hellmann--Feynman 定理，
\[
 \frac{\partial E}{\partial V_0}=-\frac{mV_0}{\hbar^2}
 =\left\langle\frac{\partial H}{\partial V_0}\right\rangle
 =-\langle\delta(x)\rangle.
\]
所以
\[
 \boxed{\langle V\rangle=-V_0\langle\delta(x)\rangle
       =-\frac{mV_0^2}{\hbar^2}=2E},\qquad
 \boxed{\langle T\rangle=E-\langle V\rangle=-E}.
\]
独立地，维里定理给出 $2\langle T\rangle=\langle xV'(x)\rangle$。
分布恒等式 $x\delta'(x)=-\delta(x)$ 可由对任意试验函数分部积分证明，故
$xV'(x)=V_0\delta(x)=-V(x)$，即 $2\langle T\rangle=-\langle V\rangle$。
与 $E=\langle T\rangle+\langle V\rangle$ 联立，仍得同一结果。
\end{enumerate}
}

\item 
二维各向同性谐振子系统哈密顿算符为$\mathrm{H}_{0}=-\frac{\hbar^{2}}{2 m}\left(\frac{\partial^{2}}{\partial x^{2}}+\frac{\partial^{2}}{\partial y^{2}}\right)+\frac{1}{2} m \omega^{2}\left(x^{2}+y^{2}\right)$。
设谐振子又受到微扰作用，总哈密顿量变成 $\mathrm{H}=H_{0}+H^{\prime}, H^{\prime}=\lambda\left(x y+y^{2}\right) $（$\lambda$ 是小常数）。
\begin{enumerate}
	%\renewcommand{\labelenumi}{\arabic{enumi}.}
	% A(\Alph) a(\alph) I(\Roman) i(\roman) 1(\arabic)
	%设定全局标号series=example	%引用全局变量resume=example
	%[topsep=-0.3em,parsep=-0.3em,itemsep=-0.3em,partopsep=-0.3em]
	%可使用leftmargin调整列表环境左边的空白长度 [leftmargin=0em]
	\item
写出 $\mathrm{H}_{0}$ 的第 $\mathrm{n}$ 级激发态能量及其简并度。
\item 
求 $\mathrm{H}_{0}$ 基态能量的一级微扰修正。

\item 
求 $\mathrm{H}_{0}$ 第一激发态能量的一级微扰修正。
	


	
\end{enumerate}
提示: 一维情况下有 $\left\langle\mathrm{n}_{\mathrm{x}}^{\prime}|\hat{x}| n_{x}\right\rangle=\sqrt{\frac{\hbar}{2 m \omega}}\left(\sqrt{n_{x}+1} \delta_{n_{x}^{\prime},n_{x}+1}+\sqrt{n_{x}} \delta_{n_{x}^{\prime},n_{x}-1}\right)$。

\banswer{
令 $x_0=\sqrt{\hbar/(2m\omega)}$。
\begin{enumerate}
\item 无扰态为 $|n_x,n_y\rangle$，其能量
\[
 \boxed{E_N^{(0)}=(N+1)\hbar\omega},\quad N=n_x+n_y.
\]
对固定 $N$，$n_x$ 可取 $0,\ldots,N$，而 $n_y=N-n_x$，所以简并度为 $N+1$。
\item 基态 $|00\rangle$ 非简并。因为 $\langle n|x|n\rangle=0$、
$\langle n|y^2|n\rangle=x_0^2(2n+1)$，
\[
 E_0^{(1)}=\lambda(\langle0|x|0\rangle\langle0|y|0\rangle
                            +\langle0|y^2|0\rangle)=\lambda x_0^2.
\]
故 $\boxed{E_0\simeq\hbar\omega+\lambda\hbar/(2m\omega)}$。
\item 固定简并基顺序为 $|0,1\rangle,|1,0\rangle$。
$xy$ 的对角元为零，而非对角元为
$\langle0|x|1\rangle\langle1|y|0\rangle=x_0^2$。
$y^2$ 的对角元分别为 $3x_0^2,x_0^2$，非对角元因 $x$ 方向正交而为零，故
\[
 W=\lambda x_0^2\begin{pmatrix}3&1\\1&1\end{pmatrix}.
\]
令 $d=E^{(1)}/(\lambda x_0^2)$，久期方程为
$(3-d)(1-d)-1=d^2-4d+2=0$，得 $d=2\pm\sqrt2$。因此
\[
 \boxed{E_{1,\pm}\simeq2\hbar\omega+
                 \frac{\lambda\hbar}{2m\omega}(2\pm\sqrt2)}.
\]
相应归一化零级态可写为
\[
 |\pm\rangle=\frac{|0,1\rangle+(-1\pm\sqrt2)|1,0\rangle}
                  {\sqrt{1+(-1\pm\sqrt2)^2}}.
\]
\end{enumerate}
}

\newpage
\item 
已知一粒子自旋为 $1 / 2$ 、磁矩为 $\mu$ 的粒子置于 $\mathrm{z}$ 方向的恒定磁场
$\mathrm{B}=B_{0}=\left(0,0, \mathrm{B}_{0}\right)$ 中。系统哈密顿算符 $\mathrm{H_{0}}=-\mu B \cdot \sigma$, 其中 $\sigma$ 为泡利矩阵。
$t<0$ 时粒子处于基态。 设系统再受到 $x$ 方向的微小脉冲磁场 $\mathrm{B}_{1}=\left(\mathrm{B}_{1}, 0,0\right) \delta(\mathrm{t})$ 作用。请
使用微扰（一级近似）进行计算:
\begin{enumerate}
	%\renewcommand{\labelenumi}{\arabic{enumi}.}
	% A(\Alph) a(\alph) I(\Roman) i(\roman) 1(\arabic)
	%设定全局标号series=example	%引用全局变量resume=example
	%[topsep=-0.3em,parsep=-0.3em,itemsep=-0.3em,partopsep=-0.3em]
	%可使用leftmargin调整列表环境左边的空白长度 [leftmargin=0em]
\item
粒子 $t>0$ 时处于激发态的概率和仍处于基态的概率。
\item
粒子 $t>0$ 时的波函数。
\item
粒子 $t>0$ 时分别处于 $\sigma_{x}$ 的两个本征态的概率（精确到 $B_{1}$ 的平方阶）。

	
\end{enumerate}


\banswer{
\begin{enumerate}
	%\renewcommand{\labelenumi}{\arabic{enumi}.}
	% A(\Alph) a(\alph) I(\Roman) i(\roman) 1(\arabic)
	%设定全局标号series=example	%引用全局变量resume=example
	%[topsep=-0.3em,parsep=-0.3em,itemsep=-0.3em,partopsep=-0.3em]
	%可使用leftmargin调整列表环境左边的空白长度 [leftmargin=0em]
	\item
		$ \sigma_{z} $表象：基态$ \ket{\psi_1}=\commamatrix[b]{1; 0} $，$ E_{1}=-\mu B_{0} $；激发态$ \ket{\psi_2}=\commamatrix[b]{0; 1} $，$ E_{2}=+\mu B_{0} $。\\
	取 $\mu B_0>0$，令 $\beta=\mu B_1/\hbar$。脉冲项为
	$H'(t)=-\mu B_1\sigma_x\delta(t)$。跨过完整脉冲的演化算符为
	\[
	 U_{\rm kick}=e^{i\beta\sigma_x}=I\cos\beta+i\sigma_x\sin\beta,
	\]
	故脉冲后 $|\psi(0^+)\rangle=\cos\beta|\psi_1\rangle+i\sin\beta|\psi_2\rangle$。
	也可由一阶含时微扰积分直接得到 $c_2^{(1)}=i\beta$。
	所以在弱脉冲 $|\beta|\ll1$ 下，
	\[
	 P_{12}=\sin^2\beta=\beta^2+O(\beta^4),\qquad
	 P_{11}=\cos^2\beta=1-\beta^2+O(\beta^4).
	\]
	这里的 $\delta(t)$ 取从脉冲前到脉冲后的完整冲量，而非端点积分的一半。
	\item 
	波函数
	\[ 
	\ket{\psi(t)} = \sqrt{1-\left(\frac{\mu B_{1}}{\hbar}\right)^{2}} \cdot  e^{-i\frac{E_{1}}{\hbar} t} \ket{\psi_1} +i\frac{\mu B_{1}}{\hbar} e^{-i\frac{E_{2}}{\hbar} t} \ket{\psi_2}
	 \]
\item
	令 $\beta=\mu B_1/\hbar$。$\sigma_x$ 的两个本征态为
	$|\pm x\rangle=(|\psi_1\rangle\pm|\psi_2\rangle)/\sqrt2$。利用上一问的波函数，得
	\[
	\begin{aligned}
	P_{+x}(t)
	&=\frac{1}{2}\left|
	\sqrt{1-\beta^2}\,e^{-iE_1t/\hbar}
	+i\beta e^{-iE_2t/\hbar}\right|^2\\
	&=\frac{1}{2}\left[1+2\beta\sqrt{1-\beta^2}
	\sin\!\left(\frac{E_2-E_1}{\hbar}t\right)\right]\\
	&=\frac{1}{2}+\beta\sin\!\left(\frac{2\mu B_0}{\hbar}t\right)+O(\beta^3),\\
	P_{-x}(t)
	&=\frac{1}{2}-\beta\sin\!\left(\frac{2\mu B_0}{\hbar}t\right)+O(\beta^3).
	\end{aligned}
	\]
	二阶项在归一化后的两个概率中恰好相消；上式因而已经精确到 $B_1^2$。
 
\end{enumerate}

	
}

\item 
考虑原子中处于同一个单电子能级 $E_{n l}$ 上的两个电子, 以 $l_{1} $、$ s_{1}$ 及
$l_{2}$、$ s_{2}$ 分别表示电子$  1 $、$ 2 $ 的轨道角动量算符和自旋角动量算符, 则总轨道角动量 算符 $\mathrm{L}=l_{1}+l_{2}$, 总自旋算符 $\mathrm{S}=s_{1}+s_{2}$。
\begin{enumerate}
	%\renewcommand{\labelenumi}{\arabic{enumi}.}
	% A(\Alph) a(\alph) I(\Roman) i(\roman) 1(\arabic)
	%设定全局标号series=example	%引用全局变量resume=example
	%[topsep=-0.3em,parsep=-0.3em,itemsep=-0.3em,partopsep=-0.3em]
	%可使用leftmargin调整列表环境左边的空白长度 [leftmargin=0em]
	\item
	写出总轨道角动量量子数 $\mathrm{L}$ 的可取值。
	\item 
	写出总自旋角动量量子数 $S$ 的可取值, 并求出 $\mathrm{S}^{2}$ 和 $\mathrm{S}_{z}$ 的共同本征态。
	\item 
	根据两个电子的全同性，讨论 $\mathrm{L}+\mathrm{S}$ 的奇偶性。
	
	
	
\end{enumerate}

\banswer{
\begin{enumerate}
	%\renewcommand{\labelenumi}{\arabic{enumi}.}
	% A(\Alph) a(\alph) I(\Roman) i(\roman) 1(\arabic)
	%设定全局标号series=example	%引用全局变量resume=example
	%[topsep=-0.3em,parsep=-0.3em,itemsep=-0.3em,partopsep=-0.3em]
	%可使用leftmargin调整列表环境左边的空白长度 [leftmargin=0em]
	\item
	$\mathrm{L}=l_{1}+l_{2}$，$ L $可能的取值为$L=0,1,2, \ldots, 2 l$
	\item 
	$S=s_{1}+s_{2}$，$ S $可能的取值为$ 0,1 $。
	$\mathrm{S}^{2}$ 和 $\mathrm{S}_{z}$ 的共同本征态$ \ket{sm} $有四个：
	\[ 
	\begin{tblr}{ll}
		|1,1\rangle&=|\uparrow \uparrow\rangle\\
		|1,0\rangle&=\frac{1}{\sqrt{2}}(|\uparrow \downarrow\rangle+|\downarrow \uparrow\rangle)\\
		|1,-1\rangle&=|\downarrow \downarrow\rangle\\
	|0,0\rangle &=\frac{1}{\sqrt{2}}(|\uparrow \downarrow\rangle-|\downarrow \uparrow\rangle) 
	\end{tblr}
	 \]
\item 
利用$ C-G $系数关于交换的关系式
\[ 
\braket{j_{1}m_{1}j_{2}m_{2}|j_{3}m_{3}}=(-1)^{j_{1}+j_{2}-j_{3}}\braket{j_{2}m_{2}j_{1}m_{1}|j_{3}m_{3}}
 \]	
	由于 $l_1=l_2=l$，轨道耦合态在交换两个电子时取得因子
	$(-1)^{2l-L}=(-1)^L$：$L$ 为偶数时空间部分对称，$L$ 为奇数时反对称。自旋三重态 $S=1$ 对称，自旋单态 $S=0$ 反对称，其交换因子为 $(-1)^{S+1}$。电子总波函数必须反对称，因此
	\[
	(-1)^L(-1)^{S+1}=-1
	\quad\Longrightarrow\quad
	(-1)^{L+S}=1.
	\]
	故允许态必须满足 $L+S$ 为偶数。
\end{enumerate}

	
}



\end{enumerate}
